\documentclass[12pt,a4paper]{article}
\usepackage[utf-8]{inputenc}
\usepackage[margin=1in]{geometry}
\usepackage{graphicx}
\usepackage{amsmath}
\usepackage{amssymb}
\usepackage{xcolor}
\usepackage{listings}
\usepackage{hyperref}
\usepackage{fancyhdr}
\usepackage{tcolorbox}
\usepackage{enumerate}
\usepackage{booktabs}
\usepackage{array}
\usepackage{textcomp}

% Configure listings for SQL code
\lstset{
    language=SQL,
    basicstyle=\ttfamily\footnotesize,
    keywordstyle=\color{blue}\bfseries,
    commentstyle=\color{gray},
    stringstyle=\color{red},
    breaklines=true,
    showstringspaces=false,
    tabsize=4,
    frame=lines,
    rulecolor=\color{black},
    numbers=left,
    numberstyle=\tiny,
    stepnumber=1
}

% Header and footer
\pagestyle{fancy}
\fancyhf{}
\lhead{MEDAITHON-04 Project}
\rhead{Interview Preparation}
\cfoot{\thepage}

\title{\textbf{\Large MEDAITHON-04 Smart Hospital Management}\\
\textbf{Full-Stack Healthcare Management System}\\
\normalsize Interview Guide \& SQL Developer Role}
\author{SQL Developer \& Database Specialist}
\date{\today}

\begin{document}

\maketitle

\tableofcontents
\newpage

% ============================================================================
\section{Project Overview}
% ============================================================================

\subsection{What is MEDAITHON-04?}

MEDAITHON-04 is a \textbf{Smart Hospital Management System} built during a 50-hour hackathon. It is a full-stack web application designed to:

\begin{itemize}
    \item Manage \textbf{patient records} and registration
    \item Schedule \textbf{medical appointments} efficiently
    \item Store \textbf{digital patient records} securely
    \item Handle \textbf{follow-up consultations} and diagnostics
    \item Track \textbf{callback requests} for patient communication
    \item Use AI for \textbf{aorta disease detection} (Aorta AI trained model)
\end{itemize}

\subsection{Real-World Application}

\begin{tcolorbox}[colback=blue!10,colframe=blue!50!black,title=Healthcare Use Case]
Hospital receives patient $\rightarrow$ Registers in system $\rightarrow$ AI analyzes health data $\rightarrow$ Schedules follow-ups $\rightarrow$ Manages appointments \textbf{automatically}
\end{tcolorbox}

\subsection{Project Technology Stack}

\begin{table}[h]
\centering
\begin{tabular}{|l|l|l|}
\hline
\textbf{Component} & \textbf{Technology} & \textbf{Purpose} \\
\hline
Frontend & HTML, CSS, JavaScript (57.8 percent) & User Interface \\
Frontend & TypeScript (34.5 percent) & Type-Safe JavaScript \\
Backend & Node.js, Express & Server Application \\
Database & MySQL / MongoDB & Data Storage \\
AI Model & Python (Trained Model) & Disease Detection \\
\hline
\end{tabular}
\end{table}

% ============================================================================
\section{Your Role: SQL Developer and Database Specialist}
% ============================================================================

\subsection{Responsibilities}

Your primary responsibilities are:

\begin{enumerate}
    \item \textbf{Design and implement database schema} (tables, relationships)
    \item \textbf{Write SQL queries} for data retrieval and manipulation
    \item \textbf{Optimize database performance} (indexing, query optimization)
    \item \textbf{Ensure data integrity} and consistency
    \item \textbf{Manage database transactions} and constraints
    \item \textbf{Create views and procedures} for complex queries
    \item \textbf{Handle database maintenance} and backups
\end{enumerate}

\subsection{Tools You Use}

\begin{tcolorbox}[colback=yellow!10,colframe=orange!50!black,title=Your Toolkit]
\begin{itemize}
    \item \textbf{MySQL}: Relational Database Management System
    \item \textbf{SQL}: Database Query Language
    \item \textbf{MySQL Workbench}: Database Design and Management
    \item \textbf{Command Line}: MySQL CLI for administration
\end{itemize}
\end{tcolorbox}

% ============================================================================
\section{Database Schema (The Heart of Your Work)}
% ============================================================================

\subsection{Overview}

The MEDAITHON-04 system uses three main database tables:

\begin{tcolorbox}[colback=cyan!10,colframe=cyan!50!black,title=Database Tables]
\begin{enumerate}
    \item \textbf{callback\_requests} - Store patient callback requests
    \item \textbf{patients} - Store patient registration information
    \item \textbf{patient\_followups} - Store follow-up consultation data
\end{enumerate}
\end{tcolorbox}

\subsection{Table 1: callback\_requests}

\textbf{Purpose:} Store callback requests from patients who want the hospital to contact them

\begin{lstlisting}[caption={callback\_requests Table Schema}]
CREATE TABLE IF NOT EXISTS callback_requests (
  id INT AUTO_INCREMENT PRIMARY KEY,
  name VARCHAR(120) NOT NULL,
  mobile VARCHAR(20) NOT NULL,
  otp VARCHAR(20) NOT NULL,
  created_at TIMESTAMP DEFAULT CURRENT_TIMESTAMP
);
\end{lstlisting}

\textbf{Columns:}
\begin{table}[h]
\centering
\begin{tabular}{|l|l|l|l|}
\hline
\textbf{Column} & \textbf{Type} & \textbf{Constraints} & \textbf{Purpose} \\
\hline
id & INT & AUTO INCREMENT, PRIMARY KEY & Unique identifier \\
name & VARCHAR(120) & NOT NULL & Patient name \\
mobile & VARCHAR(20) & NOT NULL & Phone number \\
otp & VARCHAR(20) & NOT NULL & One-time password \\
created\_at & TIMESTAMP & DEFAULT CURRENT\_TIMESTAMP & Record creation time \\
\hline
\end{tabular}
\end{table}

\subsection{Table 2: patients}

\textbf{Purpose:} Store main patient registration data

\begin{lstlisting}[caption={patients Table Schema}]
CREATE TABLE IF NOT EXISTS patients (
  id INT AUTO_INCREMENT PRIMARY KEY,
  uhid VARCHAR(50) NOT NULL UNIQUE,
  created_at TIMESTAMP DEFAULT CURRENT_TIMESTAMP
);
\end{lstlisting}

\textbf{Columns:}
\begin{table}[h]
\centering
\begin{tabular}{|l|l|l|l|}
\hline
\textbf{Column} & \textbf{Type} & \textbf{Constraints} & \textbf{Purpose} \\
\hline
id & INT & AUTO INCREMENT, PRIMARY KEY & Unique identifier \\
uhid & VARCHAR(50) & NOT NULL, UNIQUE & Universal Health ID \\
created\_at & TIMESTAMP & DEFAULT CURRENT\_TIMESTAMP & Registration time \\
\hline
\end{tabular}
\end{table}

\textbf{Note:} UHID (Universal Health ID) is a unique identifier assigned to each patient across the healthcare system.

\subsection{Table 3: patient\_followups}

\textbf{Purpose:} Store follow-up consultation records and diagnosis information

\begin{lstlisting}[caption={patient\_followups Table Schema}]
CREATE TABLE IF NOT EXISTS patient_followups (
  id INT AUTO_INCREMENT PRIMARY KEY,
  uhid VARCHAR(50) NOT NULL,
  diagnosis TEXT NOT NULL,
  severity VARCHAR(50) NOT NULL,
  next_follow_up VARCHAR(120) NOT NULL,
  created_at TIMESTAMP DEFAULT CURRENT_TIMESTAMP
);
\end{lstlisting}

\textbf{Columns:}
\begin{table}[h]
\centering
\begin{tabular}{|l|l|l|l|}
\hline
\textbf{Column} & \textbf{Type} & \textbf{Constraints} & \textbf{Purpose} \\
\hline
id & INT & AUTO INCREMENT, PRIMARY KEY & Unique identifier \\
uhid & VARCHAR(50) & NOT NULL & Links to patient \\
diagnosis & TEXT & NOT NULL & Medical diagnosis \\
severity & VARCHAR(50) & NOT NULL & Severity level \\
next\_follow\_up & VARCHAR(120) & NOT NULL & Next appointment date \\
created\_at & TIMESTAMP & DEFAULT CURRENT\_TIMESTAMP & Record creation time \\
\hline
\end{tabular}
\end{table}

% ============================================================================
\section{SQL Queries Used in Project}
% ============================================================================

\subsection{Query 1: Retrieve All Callback Requests}

\textbf{Purpose:} Fetch all callback requests sorted by creation date

\textbf{Use Case:} Hospital staff views pending callback requests

\begin{lstlisting}[caption={Get All Callback Requests}]
SELECT id, name, mobile, otp, created_at
FROM callback_requests
ORDER BY created_at DESC;
\end{lstlisting}

\textbf{Explanation:}
\begin{itemize}
    \item \textbf{SELECT}: Retrieves id, name, mobile, otp, and created\_at columns
    \item \textbf{FROM}: Source table is callback\_requests
    \item \textbf{ORDER BY}: Sorts by created\_at in descending order (newest first)
    \item \textbf{DESC}: Descending order (most recent first)
\end{itemize}

---

\subsection{Query 2: Retrieve All Patient Followups}

\textbf{Purpose:} Fetch all patient follow-up records

\textbf{Use Case:} Doctor views all patient follow-ups for the day

\begin{lstlisting}[caption={Get All Patient Followups}]
SELECT id, uhid, diagnosis, severity, next_follow_up, created_at
FROM patient_followups
ORDER BY created_at DESC;
\end{lstlisting}

---

\subsection{Query 3: Count Patients by Severity}

\textbf{Purpose:} Get statistical data on patient severity levels

\textbf{Use Case:} Hospital administration wants severity breakdown

\begin{lstlisting}[caption={Patient Count by Severity}]
SELECT severity, COUNT(*) AS total_patients
FROM patient_followups
GROUP BY severity
ORDER BY total_patients DESC;
\end{lstlisting}

\textbf{Explanation:}
\begin{itemize}
    \item \textbf{GROUP BY}: Groups records by severity level
    \item \textbf{COUNT(*)}: Counts total records in each group
    \item \textbf{AS total\_patients}: Alias for the count column
    \item \textbf{ORDER BY}: Sorts by count in descending order
\end{itemize}

\textbf{Example Output:}
\begin{table}[h]
\centering
\begin{tabular}{|l|l|}
\hline
\textbf{Severity} & \textbf{Total\_patients} \\
\hline
High & 25 \\
Medium & 18 \\
Low & 12 \\
\hline
\end{tabular}
\end{table}

---

\subsection{Query 4: Upcoming Follow-ups}

\textbf{Purpose:} Find patients with upcoming follow-up dates

\textbf{Use Case:} Schedule appointments for next week

\begin{lstlisting}[caption={Upcoming Patient Followups}]
SELECT uhid, diagnosis, next_follow_up
FROM patient_followups
WHERE next_follow_up >= CURDATE();
\end{lstlisting}

\textbf{Explanation:}
\begin{itemize}
    \item \textbf{WHERE}: Filters records based on condition
    \item \textbf{next\_follow\_up}: Column containing follow-up date
    \item \textbf{CURDATE()}: Current system date
    \item \textbf{>=}: Greater than or equal (today or later)
\end{itemize}

---

\subsection{Query 5: Join Patients with Followups}

\textbf{Purpose:} Link patient data with their follow-up records

\textbf{Use Case:} Get complete patient information with all follow-ups

\begin{lstlisting}[caption={Patient and Followup Information (JOIN)}]
SELECT p.id AS patient_id, p.uhid, pf.diagnosis, 
       pf.severity, pf.next_follow_up, pf.created_at
FROM patients p
INNER JOIN patient_followups pf ON pf.uhid = p.uhid
ORDER BY pf.created_at DESC;
\end{lstlisting}

\textbf{Explanation:}
\begin{itemize}
    \item \textbf{FROM patients p}: Source table (aliased as p)
    \item \textbf{INNER JOIN}: Combines rows from two tables
    \item \textbf{ON pf.uhid = p.uhid}: Join condition using UHID
    \item \textbf{AS patient\_id}: Alias for column
\end{itemize}

---

\subsection{Query 6: Follow-up Count by Patient and Severity}

\textbf{Purpose:} Statistical analysis of follow-ups per patient

\textbf{Use Case:} Identify high-maintenance patients

\begin{lstlisting}[caption={Followup Count Analysis}]
SELECT p.uhid, pf.severity, COUNT(*) AS followup_count
FROM patients p
INNER JOIN patient_followups pf ON pf.uhid = p.uhid
GROUP BY p.uhid, pf.severity
ORDER BY followup_count DESC;
\end{lstlisting}

\textbf{Explanation:}
\begin{itemize}
    \item \textbf{GROUP BY p.uhid, pf.severity}: Groups by patient and severity
    \item \textbf{COUNT(*)}: Counts follow-ups per group
    \item Helps identify patients needing more attention
\end{itemize}

% ============================================================================
\section{SQL Concepts Used in This Project}
% ============================================================================

\subsection{1. Data Definition Language (DDL)}

\textbf{Definition:} SQL commands to define database structure

\textbf{Commands Used:}
\begin{itemize}
    \item \textbf{CREATE DATABASE}: Create new database
    \item \textbf{CREATE TABLE}: Define table structure
    \item \textbf{ALTER TABLE}: Modify existing tables
\end{itemize}

---

\subsection{2. Data Query Language (DQL)}

\textbf{Definition:} SQL commands to retrieve data

\textbf{Commands Used:}
\begin{itemize}
    \item \textbf{SELECT}: Retrieve specific columns
    \item \textbf{FROM}: Specify table source
    \item \textbf{WHERE}: Filter rows based on conditions
    \item \textbf{ORDER BY}: Sort results
    \item \textbf{GROUP BY}: Group rows for aggregation
\end{itemize}

---

\subsection{3. Constraints}

\textbf{Definition:} Rules enforced on data

\begin{table}[h]
\centering
\begin{tabular}{|l|l|l|}
\hline
\textbf{Constraint} & \textbf{Purpose} & \textbf{Usage} \\
\hline
PRIMARY KEY & Unique identifier & id column \\
AUTO INCREMENT & Auto generate ID & Auto ID assignment \\
NOT NULL & Required field & name, mobile \\
UNIQUE & No duplicates & uhid \\
DEFAULT & Default value & created\_at timestamp \\
\hline
\end{tabular}
\end{table}

---

\subsection{4. Joins}

\textbf{Definition:} Combine data from multiple tables

\begin{itemize}
    \item \textbf{INNER JOIN}: Returns rows where both tables match
    \item Example: Link patients with their follow-ups using UHID
\end{itemize}

---

\subsection{5. Aggregate Functions}

\textbf{Definition:} Functions that perform calculations on data

\begin{table}[h]
\centering
\begin{tabular}{|l|l|l|}
\hline
\textbf{Function} & \textbf{Purpose} & \textbf{Example} \\
\hline
COUNT(*) & Count total rows & COUNT(*) patients \\
SUM() & Sum numeric values & SUM(amount) \\
AVG() & Average values & AVG(age) \\
MAX() & Maximum value & MAX(created\_at) \\
MIN() & Minimum value & MIN(created\_at) \\
\hline
\end{tabular}
\end{table}

% ============================================================================
\section{Project Pipeline: Step-by-Step Explanation}
% ============================================================================

\subsection{Step 1: Database Setup}

\subsubsection{Non-Technical Explanation}

Setting up the database is like \textbf{building the storage rooms} for the hospital:

\begin{tcolorbox}[colback=green!10,colframe=green!50!black,title=Setup Process]
\begin{enumerate}
    \item Create a database named medaithon57
    \item Create tables to store different types of information
    \item Set up rules for data (constraints)
    \item Make sure data is organized and consistent
\end{enumerate}
\end{tcolorbox}

\subsubsection{Technical Explanation}

\begin{lstlisting}[caption={Database Setup (Your Work!)}]
-- Step 1: Create database
CREATE DATABASE IF NOT EXISTS medaithon57;

-- Step 2: Use the database
USE medaithon57;

-- Step 3: Create tables with constraints
CREATE TABLE IF NOT EXISTS callback_requests (
  id INT AUTO_INCREMENT PRIMARY KEY,
  name VARCHAR(120) NOT NULL,
  mobile VARCHAR(20) NOT NULL,
  otp VARCHAR(20) NOT NULL,
  created_at TIMESTAMP DEFAULT CURRENT_TIMESTAMP
);

-- Repeat for other tables...
\end{lstlisting}

---

\subsection{Step 2: Data Insertion}

\subsubsection{Non-Technical Explanation}

After setting up tables, data flows into the system:

\begin{itemize}
    \item Patient registers using the web form
    \item Form data goes to backend (Node.js)
    \item Backend inserts data into database
    \item Data is now stored permanently
\end{itemize}

\subsubsection{Technical Explanation}

\begin{lstlisting}[caption={Inserting Data (Backend calls this)}]
-- Insert callback request
INSERT INTO callback_requests (name, mobile, otp)
VALUES ('John Doe', '9876543210', '123456');

-- Insert patient
INSERT INTO patients (uhid)
VALUES ('UH-2025-00001');

-- Insert patient follow-up
INSERT INTO patient_followups 
       (uhid, diagnosis, severity, next_follow_up)
VALUES ('UH-2025-00001', 'Aorta Disease', 'High', 
        '2025-06-15');
\end{lstlisting}

---

\subsection{Step 3: Data Retrieval (Your Main Work!)}

\subsubsection{Non-Technical Explanation}

Hospital staff needs to view patient data:

\begin{itemize}
    \item Doctor clicks ``View Patients''
    \item Frontend asks backend for data
    \item Backend runs SQL query on database
    \item Results displayed in table format
\end{itemize}

\subsubsection{Technical Explanation}

\begin{lstlisting}[caption={Retrieving Data (Your SQL Queries)}]
-- Show all patients needing follow-up today
SELECT uhid, diagnosis, severity, next_follow_up
FROM patient_followups
WHERE next_follow_up = CURDATE();

-- Show callback requests by date
SELECT name, mobile, created_at
FROM callback_requests
WHERE DATE(created_at) = CURDATE()
ORDER BY created_at ASC;
\end{lstlisting}

---

\subsection{Step 4: Data Analysis and Reporting}

\subsubsection{Non-Technical Explanation}

Hospital management needs statistics:

\begin{itemize}
    \item ``How many severe cases this week?''
    \item ``Which patients need follow-up?''
    \item ``How many callbacks completed?''
\end{itemize}

\subsubsection{Technical Explanation}

\begin{lstlisting}[caption={Analytics Queries}]
-- Cases by severity (for reporting)
SELECT severity, COUNT(*) AS count
FROM patient_followups
WHERE DATE(created_at) >= DATE_SUB(CURDATE(), 
                                   INTERVAL 7 DAY)
GROUP BY severity;

-- Patient with most follow-ups
SELECT uhid, COUNT(*) AS total_followups
FROM patient_followups
GROUP BY uhid
ORDER BY total_followups DESC
LIMIT 10;
\end{lstlisting}

% ============================================================================
\section{Interview Questions and Answers}
% ============================================================================

\subsection{Q1: What is the main purpose of this project?}

\textbf{Answer:}

MEDAITHON-04 is a \textbf{Smart Hospital Management System} built during a 50-hour hackathon to:

\begin{enumerate}
    \item Manage patient records digitally
    \item Schedule and track medical appointments
    \item Store patient follow-up information
    \item Use AI (Aorta AI) for disease detection
    \item Improve hospital efficiency and reduce manual work
\end{enumerate}

\textbf{Impact:} Hospitals can serve patients faster, reduce paperwork, and provide better care through organized data management.

---

\subsection{Q2: What is your role in this project?}

\textbf{Answer:}

My role is \textbf{SQL Developer and Database Specialist}. My responsibilities include:

\begin{enumerate}
    \item \textbf{Design database schema}: Created 3 tables (callback\_requests, patients, patient\_followups)
    \item \textbf{Write SQL queries}: For data retrieval, insertion, updating
    \item \textbf{Optimize performance}: Create indexes, optimize joins
    \item \textbf{Ensure data integrity}: Apply constraints and validations
    \item \textbf{Create reports}: Generate statistical queries for management
    \item \textbf{Maintain database}: Backups, security, updates
\end{enumerate}

\textbf{Focus:} Making sure data is stored securely, retrieved efficiently, and remains consistent.

---

\subsection{Q3: Explain the database schema you designed}

\textbf{Answer:}

The database has \textbf{3 main tables}:

\begin{enumerate}
    \item \textbf{callback\_requests}
    \begin{itemize}
        \item Stores patient callback requests
        \item Contains: id, name, mobile, otp, created\_at
        \item Purpose: Track patient contact requests
    \end{itemize}
    
    \item \textbf{patients}
    \begin{itemize}
        \item Core patient registration table
        \item Contains: id, uhid (unique ID), created\_at
        \item UHID ensures patient uniqueness across hospital
    \end{itemize}
    
    \item \textbf{patient\_followups}
    \begin{itemize}
        \item Follow-up consultation records
        \item Contains: id, uhid, diagnosis, severity, next\_follow\_up, created\_at
        \item Links to patients via UHID
    \end{itemize}
\end{enumerate}

\textbf{Relationships:} patients.uhid connects to patient\_followups.uhid (One-to-Many)

---

\subsection{Q4: Write a query to find high-severity patients}

\textbf{Answer:}

\begin{lstlisting}
SELECT p.uhid, pf.diagnosis, pf.severity, pf.next_follow_up
FROM patients p
INNER JOIN patient_followups pf ON p.uhid = pf.uhid
WHERE pf.severity = 'High'
ORDER BY pf.created_at DESC;
\end{lstlisting}

\textbf{Explanation:}
\begin{itemize}
    \item \textbf{JOIN}: Combines patient and follow-up data
    \item \textbf{WHERE pf.severity = 'High'}: Filters only high severity
    \item \textbf{ORDER BY}: Shows most recent first
    \item \textbf{Result}: All high-severity patients ready for urgent attention
\end{itemize}

---

\subsection{Q5: How would you optimize this query?}

\textbf{Answer:}

\begin{lstlisting}[caption={Optimized Query}]
-- Create index on severity for faster filtering
CREATE INDEX idx_severity ON patient_followups(severity);

-- Create index on uhid for faster joins
CREATE INDEX idx_uhid ON patient_followups(uhid);

-- Now the query runs faster
SELECT p.uhid, pf.diagnosis, pf.severity, pf.next_follow_up
FROM patients p
INNER JOIN patient_followups pf ON p.uhid = pf.uhid
WHERE pf.severity = 'High'
ORDER BY pf.created_at DESC;
\end{lstlisting}

\textbf{Optimization Techniques:}
\begin{itemize}
    \item \textbf{Indexing}: Speed up searches and joins
    \item \textbf{Select specific columns}: Don't use SELECT *
    \item \textbf{Use LIMIT}: Limit results if possible
    \item \textbf{Avoid multiple JOINs}: Use indexed foreign keys
\end{itemize}

---

\subsection{Q6: Write a query for callback statistics}

\textbf{Answer:}

\begin{lstlisting}[caption={Callback Statistics}]
SELECT 
  DATE(created_at) AS date,
  COUNT(*) AS total_callbacks,
  COUNT(DISTINCT mobile) AS unique_patients
FROM callback_requests
GROUP BY DATE(created_at)
ORDER BY date DESC;
\end{lstlisting}

\textbf{Shows:}
\begin{itemize}
    \item Callbacks per day
    \item Unique patients requesting callback
    \item Trend analysis for management
\end{itemize}

---

\subsection{Q7: How do you ensure data integrity?}

\textbf{Answer:}

\begin{enumerate}
    \item \textbf{NOT NULL constraints}: Ensure required fields
    \item \textbf{UNIQUE constraints}: Prevent duplicate UHIDs
    \item \textbf{PRIMARY KEY}: Every row uniquely identified
    \item \textbf{FOREIGN KEY}: Link related tables correctly
    \item \textbf{Data validation}: Check data types and ranges
    \item \textbf{Transactions}: Atomic operations (all or nothing)
\end{enumerate}

\begin{lstlisting}[caption={Example: Transaction}]
START TRANSACTION;

INSERT INTO patients (uhid) VALUES ('UH-2025-00002');
INSERT INTO patient_followups 
  (uhid, diagnosis, severity, next_follow_up)
VALUES ('UH-2025-00002', 'Aorta Disease', 'Medium', 
        '2025-06-20');

COMMIT; -- Both succeed or both fail
\end{lstlisting}

---

\subsection{Q8: Explain the JOIN operation used in this project}

\textbf{Answer:}

\textbf{JOIN Purpose:} Combine data from two related tables

\begin{lstlisting}[caption={JOIN Example}]
SELECT p.uhid, pf.diagnosis, pf.severity
FROM patients p
INNER JOIN patient_followups pf ON p.uhid = pf.uhid;
\end{lstlisting}

\textbf{How it works:}
\begin{table}[h]
\centering
\begin{tabular}{|l|l|l|}
\hline
\textbf{patients.uhid} & \textbf{JOIN} & \textbf{patient\_followups.uhid} \\
\hline
UH-001 & matches & UH-001 \\
UH-002 & matches & UH-002 \\
UH-003 & matches & UH-003 \\
\hline
\end{tabular}
\end{table}

\textbf{Types of JOINs:}
\begin{itemize}
    \item \textbf{INNER JOIN}: Only matching records
    \item \textbf{LEFT JOIN}: All from left table plus matches
    \item \textbf{RIGHT JOIN}: All from right table plus matches
    \item \textbf{FULL OUTER JOIN}: All records from both tables
\end{itemize}

---

\subsection{Q9: How would you handle database backups?}

\textbf{Answer:}

\begin{lstlisting}[caption={Backup and Restore}]
-- Backup database (from command line)
mysqldump -u root -p medaithon57 > backup_2025_06_15.sql

-- Restore database (from command line)
mysql -u root -p medaithon57 < backup_2025_06_15.sql

-- Backup specific table
mysqldump -u root -p medaithon57 patients > 
  patients_backup.sql
\end{lstlisting}

\textbf{Best Practices:}
\begin{itemize}
    \item Daily automated backups
    \item Store backups in secure location
    \item Test restore procedures regularly
    \item Monitor disk space
    \item Document backup schedule
\end{itemize}

---

\subsection{Q10: What data validation would you add?}

\textbf{Answer:}

\begin{lstlisting}[caption={Data Validation Checks}]
-- Check mobile number length
ALTER TABLE callback_requests
ADD CONSTRAINT check_mobile_length
CHECK (LENGTH(mobile) >= 10);

-- Check severity values
ALTER TABLE patient_followups
ADD CONSTRAINT check_severity
CHECK (severity IN ('Low', 'Medium', 'High', 'Critical'));

-- Prevent future dates for created_at
ALTER TABLE patients
ADD CONSTRAINT check_date
CHECK (created_at <= CURRENT_TIMESTAMP);
\end{lstlisting}

\textbf{Validation Types:}
\begin{itemize}
    \item \textbf{Length checks}: Mobile number minimum digits
    \item \textbf{Enumeration}: Valid severity levels only
    \item \textbf{Date checks}: No future timestamps
    \item \textbf{Type checks}: Correct data types
    \item \textbf{Range checks}: Values within acceptable range
\end{itemize}

% ============================================================================
\section{SQL Commands Reference}
% ============================================================================

\subsection{Data Definition Language (DDL)}

\begin{lstlisting}[caption={DDL Commands}]
-- Create database
CREATE DATABASE medaithon57;

-- Create table with constraints
CREATE TABLE patients (
  id INT AUTO_INCREMENT PRIMARY KEY,
  uhid VARCHAR(50) NOT NULL UNIQUE,
  created_at TIMESTAMP DEFAULT CURRENT_TIMESTAMP
);

-- Modify table
ALTER TABLE patients ADD COLUMN email VARCHAR(100);

-- Delete table
DROP TABLE patients;
\end{lstlisting}

---

\subsection{Data Query Language (DQL)}

\begin{lstlisting}[caption={DQL Commands}]
-- Basic SELECT
SELECT * FROM patients;

-- Select specific columns
SELECT uhid, created_at FROM patients;

-- Filter with WHERE
SELECT * FROM patients WHERE id = 1;

-- Multiple conditions
SELECT * FROM patients 
WHERE created_at > DATE_SUB(CURDATE(), INTERVAL 7 DAY)
AND status = 'Active';

-- ORDER BY
SELECT * FROM patients ORDER BY created_at DESC;

-- LIMIT
SELECT * FROM patients LIMIT 10;
\end{lstlisting}

---

\subsection{Aggregate Functions}

\begin{lstlisting}[caption={Aggregation}]
-- Count records
SELECT COUNT(*) FROM patients;

-- Group and count
SELECT severity, COUNT(*) 
FROM patient_followups 
GROUP BY severity;

-- Having clause (filter groups)
SELECT severity, COUNT(*) AS cnt
FROM patient_followups 
GROUP BY severity
HAVING cnt > 5;

-- Multiple aggregates
SELECT 
  severity,
  COUNT(*) AS count,
  MAX(created_at) AS latest,
  MIN(created_at) AS earliest
FROM patient_followups
GROUP BY severity;
\end{lstlisting}

% ============================================================================
\section{Common Mistakes and Solutions}
% ============================================================================

\subsection{Mistake 1: SQL Injection Attack}

\textbf{Problem:}
\begin{lstlisting}
-- UNSAFE! User input directly in query
SELECT * FROM patients 
WHERE uhid = '` + userInput + `';
\end{lstlisting}

\textbf{Solution:}
\begin{lstlisting}
-- SAFE! Use prepared statements
PREPARE stmt FROM 'SELECT * FROM patients WHERE uhid = ?';
EXECUTE stmt USING @uhid;
\end{lstlisting}

---

\subsection{Mistake 2: Missing Indexes}

\textbf{Problem:} Slow queries on large tables

\textbf{Solution:}
\begin{lstlisting}
-- Add indexes on frequently searched columns
CREATE INDEX idx_uhid ON patient_followups(uhid);
CREATE INDEX idx_severity ON patient_followups(severity);
CREATE INDEX idx_date ON patient_followups(created_at);
\end{lstlisting}

---

\subsection{Mistake 3: SELECT * Wildcard}

\textbf{Problem:} Slow performance, unnecessary data

\textbf{Solution:}
\begin{lstlisting}
-- Specify only needed columns
SELECT uhid, diagnosis, severity 
FROM patient_followups;
\end{lstlisting}

---

\subsection{Mistake 4: Lack of Error Handling}

\textbf{Problem:} Silent failures, data inconsistency

\textbf{Solution:}
\begin{lstlisting}
-- Use transactions and error handling
START TRANSACTION;

INSERT INTO patients (uhid) VALUES ('UH-2025-00003');

IF @@error_count = 0 THEN
  COMMIT;
ELSE
  ROLLBACK;
END IF;
\end{lstlisting}

% ============================================================================
\section{Summary and Key Takeaways}
% ============================================================================

\subsection{Project Overview}
\begin{itemize}
    \item \textbf{Goal:} Smart Hospital Management System
    \item \textbf{Built in:} 50-hour hackathon
    \item \textbf{Tech Stack:} Node.js backend, MySQL database, React frontend
    \item \textbf{AI Feature:} Aorta AI trained model for disease detection
\end{itemize}

\subsection{Your Role}
\begin{itemize}
    \item \textbf{SQL Developer}: Design and manage database
    \item \textbf{Query Writer}: Create efficient SQL queries
    \item \textbf{Optimizer}: Ensure fast, reliable data access
    \item \textbf{Data Custodian}: Ensure data integrity and security
\end{itemize}

\subsection{Database Highlights}
\begin{enumerate}
    \item \textbf{3 Main Tables}: callback\_requests, patients, patient\_followups
    \item \textbf{Primary Key}: Auto-increment ID for uniqueness
    \item \textbf{Foreign Relationship}: UHID links patients to follow-ups
    \item \textbf{Constraints}: NOT NULL, UNIQUE, DEFAULT values
\end{enumerate}

\subsection{Key SQL Queries}
\begin{itemize}
    \item \textbf{Retrieval}: SELECT with WHERE, ORDER BY
    \item \textbf{Analysis}: GROUP BY with COUNT aggregation
    \item \textbf{Joins}: INNER JOIN to link related tables
    \item \textbf{Filtering}: Multiple conditions for specific data
\end{itemize}

\subsection{Performance Tips}
\begin{itemize}
    \item Create indexes on frequently searched columns
    \item Use specific column names instead of SELECT *
    \item Use LIMIT to reduce returned rows
    \item Write efficient JOIN queries
    \item Monitor query execution time
\end{itemize}

% ============================================================================
\end{document}
